\section{Conclusion and open problems}
\label{sec:conclusion}

The four models isolate two effects of a unit preliminary action: it may
reduce the work needed to finish a job, and it may reveal where that job
belongs in the schedule.  Their exact phase diagrams make the separation
quantitative.  With revelation, both optimal competitive ratios remain
bounded as the raw cap $u$ grows and converge to the obligatory-testing
constants.  Without revelation, the blind-optimization ratios grow as
$\Theta(\sqrt u)$---deterministically with leading constant one and under an
oblivious adversary with leading constant one half.  Thus advance
information changes not merely a constant but the asymptotic dependence on
the available speedup.

The reusable part of the analysis is the completion-density description of
the fluid optimum.  It turns a distribution of hidden lengths into a
pointwise optimal completion curve, and its finite-to-fluid transfer shows
that a universal online algorithm can learn and implement that curve from a
sublinear private sample.  The resulting guarantee is instance-specific in
a precise sense: it is in expectation, compares against a multiset-aware but
placement-oblivious online competitor, and matches the leading $n^2$ cost up
to $o(n^2)$.  In blind execution this additive statement need not imply a
multiplicative bound when the clairvoyant optimum itself is $o(n^2)$.

Randomization plays two related roles.  A private shuffle converts the
unseen labels into a representative without-replacement population, and a
pilot sample then supports a stationary threshold or block plan.  The
deterministic lower bounds show what fails without that stationarity: an
adversary can correlate the revealed prefix with the algorithm's history,
forcing adaptive thresholds and, in some models, ruling out deterministic
instance optimality altogether.  Boundedness is equally structural for the
learned per-instance guarantees; rare unbounded values cannot be certified
from a sublinear pilot without a fallback or an additional loss.

Several directions remain open.  Nonuniform test costs would replace the
single completion-density module by a heterogeneous selection problem.
Weights and multiple machines would break the present one-dimensional
shortest-first envelope, while an adaptive adversary would remove the clean
private-shuffle reduction.  It would also be useful to complement the
size-asymptotic theory with explicit finite-$n$ algorithms and tuning rules,
and to determine whether sharper rates can be obtained without the slowly
growing grids and cutoffs used here.  These questions offer natural tests of
how far completion-density methods extend beyond uniform single-machine
testing.
